\section{Discussion and Limitations}
The V40-v2 results support a narrow system-level conclusion: under the fixed V40-v2 component-disjoint validation protocol, the pre-specified reliability-aware $p=0.15$ system outperforms matched early fusion in the two fixed training runs and under deterministic removal of a single modality group. The comparison does not isolate the causal contribution of the gate alone because the reliability model is trained with modality dropout while early fusion is not. The absence of trained RGB-only, thermal-only, event-only, and static-global-weight controls also limits what can be concluded about individual modalities or the specific mechanism of input-conditioned weighting.

The split protocol is stronger than an exact-duplicate check but has a defined scope. It enforces component disjointness for the locked graph composed of exact RGB content, perceptual-hash candidates, and human-adjudicated adjacent-or-near-identical observations. It cannot demonstrate that all possible scene correlations are absent. Filename identifiers were not treated as verified temporal metadata. The guard partition is not a test set and is not used to make any performance claim.

The experimental evidence is validation-only. The V40 validation partition participates in within-run checkpoint retention and final reporting. A separately acquired and locked tri-modal test set would be required for independent-test or external-generalization claims. Synthetic channel removal is a controlled zeroing operation, not an estimate of real physical sensor reliability.

\section{Conclusion}
RA-RepDet is a lightweight reliability-aware RGB-thermal-event detector evaluated against matched early fusion on a V40-v2 expanded-adjacency component-disjoint validation partition. The pre-specified reliability-aware $p=0.15$ system improves two-run mean AP50, AP75, F1, precision, and recall, and its advantage remains visible under deterministic synthetic removal of RGB, thermal, or event inputs. The model carries modest parameter and compute overhead but higher peak memory and slightly higher measured latency. The contribution is therefore a validation-only, reproducible fusion result with a transparent partition-audit scope, rather than a claim of independent testing, universal leakage absence, or physical sensor-failure robustness.
